\documentclass[12pt,a4paper]{ctexbook}
\usepackage{geometry}
\geometry{margin=2.2cm}
\usepackage{setspace}
\setstretch{1.35}
\usepackage{hyperref}
\title{全唐诗选·poet.tang.5000}
\author{古典文献整理}
\date{}
\begin{document}
\maketitle
\tableofcontents
\newpage
\section{夜泊越州逢北使}
\textbf{沈佺期}\par\vspace{0.5em}
天地降雷雨，放逐還國都。\par
重以風潮事，年月戒回艫。\par
容顏荒外老，心想域中愚。\par
憩泊在茲夜，炎雲逐斗樞。\par
䬓\{風兪\}縈海若，霹靂耿天吳。\par
鼇抃羣島失，鯨吞衆流輸。\par
偶逢金華使，握手淚相濡。\par
飢共噬齊棗，眠共席秦蒲。\par
既北思攸濟，將南睿所圖。\par
往來固無咎，何忽憚前桴。\par

\section{紹隆寺}
\textbf{沈佺期}\par\vspace{0.5em}
吾從釋迦久，無上師涅槃。\par
探道三十載，得道天南端。\par
非勝適殊方，起諠歸理難。\par
放棄乃良緣，世慮不曾干。\par
香界縈北渚，花龕隱南巒。\par
危昂階下石，演漾窗中瀾。\par
雲蓋看木秀，天空見藤盤。\par
處俗勒宴坐，居貧業行壇。\par
試將有漏軀，聊作無生觀。\par
了然究諸品，彌覺靜者安。\par

\section{神龍初廢逐南荒途出郴口北望蘇耽山}
\textbf{沈佺期}\par\vspace{0.5em}
少曾讀仙史，知有蘇耽君。\par
流望來南國，依然會昔聞。\par
泊舟問耆老，遙指孤山雲。\par
孤山郴郡北，不與衆山羣。\par
重崖下縈映，嶛嶢上乣紛。\par
碧峰泉附落，紅壁樹傍分。\par
選地今方爾，昇天因可云。\par
不才予竄迹，羽化子遺芬。\par
將覽成麟鳳，旋驚禦鬼文。\par
此中迷出處，含思獨氛氳。\par

\section{初達驩州}
\textbf{沈佺期}\par\vspace{0.5em}
流子一十八，命予偏不偶。\par
配遠天遂窮，到遲日最後。\par
水行儋耳國，陸行雕題藪。\par
魂魄遊鬼門，骸骨遺鯨口。\par
夜則忍飢臥，朝則抱病走。\par
搔首向南荒，拭淚看北斗。\par
何年赦書來？重飲洛陽酒。\par

\section{被彈}
\textbf{沈佺期}\par\vspace{0.5em}
知人昔不易，舉非貴易失。\par
爾何按國章，無罪見呵叱。\par
平生守直道，遂爲衆所嫉。\par
少以文作吏，手不曾開律。\par
一旦法相持，荒忙意如漆。\par
幼子雙囹圄，老夫一念室。\par
昆弟兩三人，相次俱囚桎。\par
萬鑠當衆怒，千謗無片實。\par
庶以白黑讒，顯此涇渭質。\par
劾吏何咆哮，晨夜聞撲抶。\par
事間拾虛證，理外存枉筆。\par
懷痛不見伸，抱冤竟難悉。\par
窮囚多垢膩，愁坐饒蟣虱。\par
三日唯一飯，兩旬不再櫛。\par
是時盛夏中，暵赫多瘵疾。\par
瞪目眠欲閉，暗嗚氣不出。\par
有風自扶搖，鼓蕩無倫匹。\par
安得吹浮雲，令我見白日。\par

\section{枉繫二首 一}
\textbf{沈佺期}\par\vspace{0.5em}
吾憐曾家子，昔有投杼疑。\par
吾憐姬公旦，非無鴟鴞詩。\par
臣子竭忠孝，君親惑讒欺。\par
萋斐離骨肉，含愁興此辭。\par

\section{枉繫二首 二}
\textbf{沈佺期}\par\vspace{0.5em}
昔日公冶長，非罪遇縲絏。\par
聖人降其子，古來歎獨絕。\par
我無毫髪瑕，苦心懷冰雪。\par
今代多秀士，誰能繼明轍。\par

\section{黃鶴}
\textbf{沈佺期}\par\vspace{0.5em}
黃鶴佐丹鳳，不能羣白鷴。\par
拂雲遊四海，弄影到三山。\par
遙憶君軒上，來下天池間。\par
明珠世不重，知有報恩環。\par

\section{傷王學土}
\textbf{沈佺期}\par\vspace{0.5em}
閉囚斷外事，昧坐半餘朞。\par
有言頴叔子，亡來已一時。\par
初聞宛不信，中話涕漣洏。\par
痛哉玄夜重，何遽青春姿。\par
憶汝曾旅食，屢空瀍澗湄。\par
吾徒祿未厚，筲斗愧相貽。\par
原憲貧無愁，顏回樂自持。\par
詔書擇才善，君爲王子師。\par
寵儒名可尚，論秩官猶欺。\par
化往不復見，情來安可思。\par
目絕毫翰灑，耳無歌諷期。\par
靈柩寄何處，精魂今何之？\par
恨予在丹棘，不得看素旗。\par
孀妻知己歎，幼子路人悲。\par
感遊值商日，絕弦留此詞。\par

\section{古鏡}
\textbf{沈佺期}\par\vspace{0.5em}
莓苔翳清池，蝦蟆蝕明月。\par
埋落今如此，照心未嘗歇。\par
願垂拂拭恩，爲君鑒玄髪。\par

\section{鳳簫曲}
\textbf{沈佺期}\par\vspace{0.5em}
八月涼風動高閣，千金麗人捲綃幕。\par
已憐池上歇芳菲，不念君恩坐搖落。\par
世上榮華如轉蓬，朝隨阡陌暮雲中。\par
飛燕侍寢昭陽殿，班姬飲恨長信宮。\par
長信宮，昭陽殿，春來歌舞妾自知，秋至簾櫳君不見。\par
昔時嬴女厭世紛，學吹鳳簫乘彩雲。\par
含情轉睞向蕭史，千載紅顏持贈君。\par

\section{古歌}
\textbf{沈佺期}\par\vspace{0.5em}
落葉流風向玉臺，夜寒秋思洞房開。\par
水晶簾外金波下，雲母窗前銀漢回。\par
玉階陰陰苔蘚色，君王屢綦難再得。\par
璇閨窈窕秋夜長，繡戶徘徊明月光。\par
燕姬綵帳芙蓉色，秦女金爐蘭麝香。\par
北斗七星橫夜半，清歌一曲斷君腸。\par

\section{七夕曝衣篇}
\textbf{沈佺期}\par\vspace{0.5em}
君不見昔日宜春太液邊，披香畫閣與天連。\par
燈火灼爍九微暎，香氣氛氳百和然。\par
此夜星繁河正白，人傳織女牽牛客。\par
宮中擾擾曝衣樓，天上娥娥紅粉席。\par
曝衣何許曛半黃，宮中綵女提玉箱。\par
珠履奔騰上蘭砌，金梯宛轉出梅梁。\par
絳河裏，碧煙上，雙花伏兔畫屏風，四子盤龍擎斗帳。\par
舒羅散縠雲霧開，綴玉垂珠星漢回。\par
朝霞散彩羞衣架，晚月分光劣鏡臺。\par
上有仙人長命綹，中看玉女迎歡繡。\par
玳瑁簾中別作春，珊瑚窗裏翻成晝。\par
椒房金屋寵新流，意氣嬌奢不自由。\par
漢文宜惜露臺費，晉武須焚前殿裘。\par

\section{入少密溪}
\textbf{沈佺期}\par\vspace{0.5em}
雲峯苔壁繞溪斜，江路香風夾岸花。\par
樹密不言通鳥道，雞鳴始覺有人家。\par
人家更在深巖口，澗水周流宅前後。\par
遊魚瞥瞥雙釣童，伐木丁丁一樵叟。\par
自言避喧非避秦，薜衣耕鑿帝堯人。\par
相留且待雞黍熟，夕臥深山蘿月春。\par

\section{霹靂引}
\textbf{沈佺期}\par\vspace{0.5em}
歲七月，火伏而金生。\par
客有鼓琴於門者，奏霹靂之商聲。\par
始戛羽以騞砉，終扣宮而砰駖。\par
電耀耀兮龍躍，雷闐闐兮雨冥。\par
氣嗚唅以會雅，態歘翕以橫生。\par
有如驅千旗，制五兵。\par
截荒虺，斮長鯨。\par
孰與廣陵比，意別鶴儔精而已。\par
俾我雄子魄動，毅夫髪立。\par
懷恩不淺，武義雙輯。\par
視胡若芥，剪羯如拾。\par
豈徒慷慨中筵，備羣娛之翕習哉。\par

\section{立春日內出綵花應制}
\textbf{沈佺期}\par\vspace{0.5em}
合殿春應早，開箱綵預知。\par
花迎宸翰發，葉待御筵披。\par
梅訝香全少，桃驚色頓移。\par
輕生承剪拂，長伴萬年枝。\par

\section{晦日滻水應制}
\textbf{沈佺期}\par\vspace{0.5em}
素滻接宸居，青門盛祓除。\par
摘蘭喧鳳野，浮藻溢龍渠。\par
苑蝶飛殊懶，宮鶯囀不疎。\par
星移天上入，歌舞向儲胥。\par

\section{奉和洛陽翫雪應制}
\textbf{沈佺期}\par\vspace{0.5em}
周王甲子旦，漢后德陽宮。\par
灑瑞天庭裏，驚春御苑中。\par
氛氳生浩氣，颯杳舞回風。\par
宸藻光盈尺，賡歌樂歲豐。\par

\section{三日梨園侍宴}
\textbf{沈佺期}\par\vspace{0.5em}
九重馳道出，三巳禊堂開。\par
畫鷁中流動，青龍上苑來。\par
野花飄御座，河柳拂天杯。\par
日晚迎祥處，笙鏞下帝臺。\par

\section{幸梨園亭觀打毬應制}
\textbf{沈佺期}\par\vspace{0.5em}
今春芳苑遊，接武上瓊樓。\par
宛轉縈香騎，飄颻拂畫毬。\par
俯身迎未落，迴轡逐傍流。\par
秖爲看花鳥，時時誤失籌。\par

\section{九日臨渭亭侍宴應制得長字}
\textbf{沈佺期}\par\vspace{0.5em}
御氣幸金方，憑高薦羽觴。\par
魏文頒菊蕋，漢武賜萸房。\par
秋變銅池色，晴添銀樹光。\par
年年重九慶，日月奉天長。\par

\section{歲夜安樂公主滿月侍宴}
\textbf{沈佺期}\par\vspace{0.5em}
除夜子星回，天孫滿月杯。\par
詠歌麟趾合，簫管鳳雛來。\par
歲炬常然桂，春盤預折梅。\par
聖皇千萬壽，垂曉御樓開。\par

\section{安樂公主移入新宅}
\textbf{沈佺期}\par\vspace{0.5em}
初聞衡漢來，移住斗城隈。\par
錦帳迎風轉，瓊筵拂霧開。\par
馬香遺舊埒，鳳吹繞新臺。\par

\section{仙萼亭初成侍宴應制}
\textbf{沈佺期}\par\vspace{0.5em}
山中氣色和，宸賞第中過。\par
輦路披仙掌，帷宮拂帝蘿。\par
泉臨香澗落，峰入翠雲多。\par
無異登玄圃，東南望白河。\par

\section{送金城公主適西蕃應制}
\textbf{沈佺期}\par\vspace{0.5em}
金牓扶丹掖，銀河屬紫閽。\par
那堪將鳳女，還以嫁烏孫。\par
玉就歌中怨，珠辭掌上恩。\par
西戎非我匹，明主至公存。\par

\section{幸白鹿觀應制}
\textbf{沈佺期}\par\vspace{0.5em}
紫鳳真人府，斑龍太上家。\par
天流芝蓋下，山轉桂旗斜。\par
聖藻垂寒露，仙杯落晚霞。\par
唯應問王母，桃作幾時花。\par

\section{洛陽道}
\textbf{沈佺期}\par\vspace{0.5em}
九門開洛邑，雙闕對河橋。\par
白日青春道，軒裳半下朝。\par
乘羊稚子看，拾翠美人嬌。\par
行樂歸恆晚，香塵撲地遙。\par

\section{驄馬}
\textbf{沈佺期}\par\vspace{0.5em}
西北五花驄，來時道向東。\par
四蹄碧玉片，雙眼黃金瞳。\par
鞍上留明月，嘶間動朔風。\par
借君馳沛艾，一戰取雲中。\par

\section{銅雀臺}
\textbf{沈佺期}\par\vspace{0.5em}
昔年分鼎地，今日望陵臺。\par
一旦雄圖盡，千秋遺令開。\par
綺羅君不見，歌舞妾空來。\par
恩共漳河水，東流無重回。\par

\section{長門怨}
\textbf{沈佺期}\par\vspace{0.5em}
月皎風泠泠，長門次掖庭。\par
玉階聞墜葉，羅幌見飛螢。\par
清露凝珠綴，流塵下翠屏。\par
妾心君未察，愁歎劇繁星。\par

\section{巫山高二首 一}
\textbf{沈佺期}\par\vspace{0.5em}
巫山峯十二，合沓隱昭回。\par
俯眺琵琶峽，平看雲雨臺。\par
古槎天外倚，瀑水日邊來。\par
何忍猿啼夜，荆王枕席開。\par

\section{巫山高二首 二}
\textbf{沈佺期}\par\vspace{0.5em}
神女向高唐，巫山下夕陽。\par
裴回作行雨，婉孌逐荆王。\par
電影江前落，雷聲峽外長。\par
霽雲無處所，臺館曉蒼蒼。\par

\section{巫山高}
\textbf{沈佺期}\par\vspace{0.5em}
巫山高不極，合沓狀奇新。\par
暗谷疑風雨，陰崖若鬼神。\par
月明三峽曙，潮滿九江春。\par
爲問陽臺客，應知入夢人。\par

\section{七夕}
\textbf{沈佺期}\par\vspace{0.5em}
秋近雁行稀，天高鵲夜飛。\par
妝成應懶織，今夕渡河歸。\par
月皎宜穿線，風輕得曝衣。\par
來時不可覺，神驗有光輝。\par

\section{春閨}
\textbf{沈佺期}\par\vspace{0.5em}
鐵馬三軍去，金閨二月還。\par
邊愁離上國，春夢失陽關。\par
池水琉璃淨，園花玳瑁斑。\par
歲華空自擲，憂思不勝顏。\par

\section{奉和聖製同皇太子遊慈恩寺應制}
\textbf{沈佺期}\par\vspace{0.5em}
肅肅蓮花界，熒熒貝葉宮。\par
金人來夢裏，白馬出城中。\par
湧塔初從地，焚香欲遍空。\par
天歌應春籥，非是爲春風。\par

\section{和洛州康士曹庭芝望月有懷}
\textbf{沈佺期}\par\vspace{0.5em}
天使下西樓，光含萬象秋。\par
臺前疑挂鏡，簾外似懸鉤。\par
張尹將眉學，班姬取扇儔。\par
佳期應借問，爲報在刀頭。\par

\section{壽陽王花燭}
\textbf{沈佺期}\par\vspace{0.5em}
仙媛乘龍夕，天孫捧雁來。\par
可憐桃李樹，更遶鳳皇臺。\par
燭送香車入，花臨寶扇開。\par
莫令銀箭曉，爲盡合歡杯。\par

\section{隴頭水}
\textbf{沈佺期}\par\vspace{0.5em}
隴山飛落葉，隴雁度寒天。\par
愁見三秋水，分爲兩地泉。\par
西流入羗郡，東下向秦川。\par
征客重回首，肝腸空自憐。\par

\section{關山月}
\textbf{沈佺期}\par\vspace{0.5em}
漢月生遼海，朣朧出半暉。\par
合昏玄菟郡，中夜白登圍。\par
暈落關山迥，光含霜霰微。\par
將軍聽曉角，戰馬欲南歸。\par

\section{折楊柳}
\textbf{沈佺期}\par\vspace{0.5em}
玉窗朝日暎，羅帳春風吹。\par
拭淚攀楊柳，長條踠地垂。\par
白花飛歷亂，黃鳥思參差。\par
妾自肝腸斷，傍人那得知。\par

\section{梅花落}
\textbf{沈佺期}\par\vspace{0.5em}
鐵騎幾時回，金閨怨早梅。\par
雪寒花已落，風暖葉應開。\par
夕逐新春管，香迎小歲杯。\par
盛時何足貴，書裏報輪臺。\par

\section{紫騮馬}
\textbf{沈佺期}\par\vspace{0.5em}
青玉紫騮鞍，驕多影屢盤。\par
荷君能剪拂，躞蹀噴桑乾。\par
踠足追奔易，長鳴遇賞難。\par
摐金一萬里，霜露不辭寒。\par

\section{上之回}
\textbf{沈佺期}\par\vspace{0.5em}
制書下關右，天子問回中。\par
壇墠經過遠，威儀侍從雄。\par
黃麾搖晝日，青幰曳松風。\par
迴望甘泉道，龍山隱漢宮。\par

\section{王昭君}
\textbf{沈佺期}\par\vspace{0.5em}
非君惜鸞殿，非妾妬蛾眉。\par
薄命由驕虜，無情是畫師。\par
嫁來胡地日，不並漢宮時。\par
心苦無聊賴，何堪馬上辭。\par

\section{被試出塞}
\textbf{沈佺期}\par\vspace{0.5em}
十年通大漠，萬里出長平。\par
寒日生戈劒，陰雲拂斾旌。\par
飢烏啼舊壘，疲馬戀空城。\par
辛苦臯蘭北，胡霜損漢兵。\par

\section{牛女}
\textbf{沈佺期}\par\vspace{0.5em}
粉席秋期緩，針樓別怨多。\par
奔龍爭度日，飛鵲亂填河。\par
失喜先臨鏡，含羞未解羅。\par
誰能留夜色，來夕倍還梭。\par

\section{雜詩三首 一}
\textbf{沈佺期}\par\vspace{0.5em}
落葉驚秋婦，高砧促暝機。\par
蜘蛛尋月度，螢火傍人飛。\par
清鏡紅埃入，孤燈綠焰微。\par
怨啼能至曉，獨自懶縫衣。\par

\section{雜詩三首 二}
\textbf{沈佺期}\par\vspace{0.5em}
妾家臨渭北，春夢著遼西。\par
何苦朝鮮郡，年年事鼓鼙。\par
燕來紅壁語，鶯向綠窗啼。\par
爲許長相憶，闌干玉筯齊。\par

\section{雜詩三首 三}
\textbf{沈佺期}\par\vspace{0.5em}
聞道黃龍戍，頻年不解兵。\par
可憐閨裏月，長在漢家營。\par
少婦今春意，良人昨夜情。\par
誰能將旗鼓，一爲取龍城。\par

\section{剪綵}
\textbf{沈佺期}\par\vspace{0.5em}
宮女憐芳樹，裁花競早榮。\par
寒依刀尺盡，春向綺羅生。\par
弱蔕盤絲發，香蕤結素成。\par
纖枝幸不棄，長就玉階傾。\par

\section{和中書侍郎楊再思春夜宿直}
\textbf{沈佺期}\par\vspace{0.5em}
西禁青春滿，南端皓月微。\par
千廬宵駕合，五夜曉鐘稀。\par
星斗橫綸閣，天河度瑣闈。\par
煙光章奏裏，紛向夕郎飛。\par

\section{和常州崔使君寒食夜}
\textbf{沈佺期}\par\vspace{0.5em}
聞道清明近，春闈向夕闌。\par
行遊晝不厭，風物夜宜看。\par
斗柄更初轉，梅香暗裏殘。\par
無勞秉華燭，晴月在南端。\par

\section{和崔正諫登秋日早朝}
\textbf{沈佺期}\par\vspace{0.5em}
雞鳴朝謁滿，露白禁門秋。\par
爽氣臨旌戟，朝光暎冕旒。\par
河宗來獻寶，天子命焚裘。\par
獨負池陽議，言從建禮遊。\par

\section{答甯處州書}
\textbf{沈佺期}\par\vspace{0.5em}
書報天中赦，人從海上聞。\par
九泉開白日，六翮起青雲。\par
質幸恩先貸，情孤枉未分。\par
自憐涇渭別，誰與奏明君。\par

\section{李舍人山園送龐邵}
\textbf{沈佺期}\par\vspace{0.5em}
符傳有光輝，諠諠出帝畿。\par
東隣借山水，南陌駐驂騑。\par
握手涼風至，當歌秋日微。\par
高幨去勿緩，人吏待霜威。\par

\section{送陸侍御餘慶北使}
\textbf{沈佺期}\par\vspace{0.5em}
古人貴將命，之子出輶軒。\par
受委當不辱，隨時敢贈言。\par
朔途際遼海，春思繞轘轅。\par
安得回白日，留歡盡綠樽。\par

\section{洛州蕭司兵謁兄還赴洛成禮}
\textbf{沈佺期}\par\vspace{0.5em}
棠棣日光輝，高襟應序歸。\par
來成鴻雁聚，去作鳳皇飛。\par
細草乘輕傳，驚花慘別衣。\par
灞亭春有酒，岐路惜芬菲。\par

\section{餞高唐州詢}
\textbf{沈佺期}\par\vspace{0.5em}
弱冠相知早，中年不見多。\par
生涯在王事，客鬢各蹉跎。\par
良守初分岳，嘉聲即潤河。\par
還從漢闕下，傾耳聽中和。\par

\section{餞唐郎中洛陽令}
\textbf{沈佺期}\par\vspace{0.5em}
一臺推往妙，三史佇來修。\par
應宰鳧還集，辭郎雉少留。\par
郊筵乘落景，亭傳理殘秋。\par
願以弦歌暇，芝蘭想舊遊。\par

\section{樂城白鶴寺}
\textbf{沈佺期}\par\vspace{0.5em}
碧海開龍藏，青雲起雁堂。\par
潮聲迎法鼓，雨氣濕天香。\par
樹接前山暗，溪承瀑水涼。\par
無言謫居遠，清淨得空王。\par

\section{遊少林寺}
\textbf{沈佺期}\par\vspace{0.5em}
長歌遊寶地，徙倚對珠林。\par
雁塔風霜古，龍池歲月深。\par
紺園澄夕霽，碧殿下秋陰。\par
歸路煙霞晚，山蟬處處吟。\par

\section{嶽館}
\textbf{沈佺期}\par\vspace{0.5em}
洞壑仙人館，孤峰玉女臺。\par
空濛朝氣合，窈窕夕陽開。\par
流澗含輕雨，虛巖應薄雷。\par
正逢鸞與鶴，歌舞出天來。\par

\section{早發平昌島}
\textbf{沈佺期}\par\vspace{0.5em}
解纜春風後，鳴榔曉漲前。\par
陽烏出海樹，雲雁下江煙。\par
積氣衝長島，浮光溢大川。\par
不能懷魏闕，心賞獨泠然。\par

\section{夜宿七盤嶺}
\textbf{沈佺期}\par\vspace{0.5em}
獨遊千里外，高臥七盤西。\par
曉月臨窗近，天河入戶低。\par
芳春平仲綠，清夜子規啼。\par
浮客空留聽，褒城聞曙雞。\par

\section{十三四時嘗從巫峽過他日偶然有思}
\textbf{沈佺期}\par\vspace{0.5em}
小度巫山峽，荆南春欲分。\par
使君灘上草，神女館前雲。\par
樹悉江中見，猿多天外聞。\par
別來如夢裏，一想一氛氳。\par

\section{初達驩州}
\textbf{沈佺期}\par\vspace{0.5em}
自昔聞銅柱，行來向一年。\par
不知林邑地，猶隔道明天。\par
雨露何時及？京華若個邊。\par
思君無限淚，堪作日南泉。\par

\section{嶺表逢寒食}
\textbf{沈佺期}\par\vspace{0.5em}
嶺外無寒食，春來不見餳。\par
洛陽新甲子，何日是清明？\par
花柳爭朝發，軒車滿路迎。\par
帝鄉遙可念，腸斷報親情。\par

\section{驩州南亭夜望}
\textbf{沈佺期}\par\vspace{0.5em}
昨夜南亭望，分明夢洛中。\par
室家誰道別，兒女案嘗同。\par
忽覺猶言是，沈思始悟空。\par
肝腸餘幾寸，拭淚坐春風。\par

\section{少遊荆湘因有是題}
\textbf{沈佺期}\par\vspace{0.5em}
峴北焚蛟浦，巴東射雉田。\par
歲時宜楚俗，耆舊在襄川。\par
憶昨經過處，離今二十年。\par
因君訪生死，相識幾人全。\par

\section{咸陽覽古}
\textbf{沈佺期}\par\vspace{0.5em}
咸陽秦帝居，千載坐盈虛。\par
版築林光盡，壇場霤聽疎。\par
野橋疑望日，山火類焚書。\par
唯有驪峰在，空聞厚葬餘。\par

\section{覽鏡}
\textbf{沈佺期}\par\vspace{0.5em}
霏霏日搖蕙，騷騷風灑蓮。\par
時芳固相奪，俗態豈恆堅。\par
恍忽夜川裏，蹉跎朝鏡前。\par
紅顏與壯志，太息此流年。\par

\section{題椰子樹}
\textbf{沈佺期}\par\vspace{0.5em}
日南椰子樹，香褭出風塵。\par
叢生調木首，圓實檳榔身。\par
玉房九霄露，碧葉四時春。\par
不及塗林果，移根隨漢臣。\par

\section{同獄者歎獄中無燕}
\textbf{沈佺期}\par\vspace{0.5em}
何許乘春燕，多知辨夏臺。\par
三時欲併盡，雙影未嘗來。\par
食蕋嫌叢棘，銜泥怯死灰。\par
不如黃雀語，能雪冶長猜。\par

\section{則天門赦改年}
\textbf{沈佺期}\par\vspace{0.5em}
聖人宥天下，幽鑰動圜狴。\par
六甲迎黃氣，三元降紫泥。\par
籠僮上西鼓，振迅廣陽雞。\par
歌舞將金帛，汪洋被遠黎。\par

\section{喜赦}
\textbf{沈佺期}\par\vspace{0.5em}
去歲投荒客，今春肆眚歸。\par
律通幽谷暖，盆舉太陽輝。\par
喜氣迎冤氣，青衣報白衣。\par
還將合浦葉，俱向洛城飛。\par

\section{秦州薛都督挽詞}
\textbf{沈佺期}\par\vspace{0.5em}
十里絳山幽，千年汾水流。\par
碑傳門客建，劒是故人留。\par
隴樹烟含夕，山門月對秋。\par
古來鐘鼎盛，共盡一蒿丘。\par

\section{天官崔侍郎夫人盧氏挽歌}
\textbf{沈佺期}\par\vspace{0.5em}
偕老言何謬，香魂事永違。\par
潘魚從此隔，陳鳳宛然飛。\par
埋鏡泉中暗，藏鐙地下微。\par
猶憑少君術，彷彿覩容輝。\par

\section{章懷太子靖妃挽詞}
\textbf{沈佺期}\par\vspace{0.5em}
彤史佳聲載，青宮懿範留。\par
形將鸞鏡隱，魂伴鳳笙遊。\par
送馬嘶殘日，新螢落晚秋。\par
不知蒿里曙，空見隴雲愁。\par

\section{奉和立春遊苑迎春}
\textbf{沈佺期}\par\vspace{0.5em}
東郊暫轉迎春仗，上苑初飛行慶杯。\par
風射蛟冰千片斷，氣衝魚鑰九關開。\par
林中覓草纔生蕙，殿裏爭花併是梅。\par
歌吹銜恩歸路晚，棲烏半下鳳城來。\par

\end{document}